\chapter{IVPs and BVPs}

Our entry point into the study of algorithms for solving differential 
equations (DEs) is two types of ODEs,i.e., initial value problems (IVPs) and boundary value problems (BVPs). Let us describe these for an example ODE: 
\begin{equation*}
    \frac{d^2 u}{dt^2} + u = 0.
\end{equation*}
By itself this ODE is not well-posed since it may have infinitely 
many solutions. In fact, any function of the form 
\begin{equation*}
    u(t) = A \cos(t) + B \sin(t),
\end{equation*}
with arbitrary constants $A$ and $B$ 
is a solution; simply plug it into the ODE to verify. So we need additional 
side information to uniquely determine a solution. This is where 
the difference between IVPs and BVPs arises: 

\begin{itemize}
    \item In an IVP the additional data is provided at a fixed 
    point in time, typically at the initial time $t_0$. 
    For example, for our ODE, an IVP could be specified as
    \begin{equation*}
        u(0) = u_0, \quad \frac{du}{dt}(0) = v_0,
    \end{equation*}
    where $u_0$ and $v_0$ are given constants and we took $t_0 = 0$. In this case, the 
    solution is  
    \begin{equation*}
        u(t) = u_0 \cos(t) + v_0 \sin(t). 
    \end{equation*}

    \item In a BVP the additional data is provided at the "boundaries"
     of a domain, typically at two different points in time, 
     say $t_0$ and $t_1$. 
    For our ODE, a BVP could be specified as
    \begin{equation*}
        u(0) = u_0, \quad u(\pi/2) = u_1,
    \end{equation*}
    where $u_0$ and $u_1$ are given constants and we took $t_0 = 0$ 
    and $t_1 = \pi/2$. In this case, the 
    solution must satisfy the boundary conditions at 
    both $t_0$ and $t_1$ and so we obtain 
    \begin{equation*}
        u(t) = u_0 \cos(t) + u_1 \sin(t).
       \end{equation*}
\end{itemize}

We readily notice a difference between the two problems. The IVP with 
the above initial conditions (think initial location and velocity) always 
has a unique solution, i.e., we can find $A$ and $B$ via the formulae: 
\begin{equation*}
    A = u_0, \quad B = v_0.
\end{equation*}
On the other hand, the BVP may have infinitely many solutions or
may have no solutions at all.
For instance, consider the boundary data
    \begin{equation*}
        u(0) = c_1, \quad u(\pi) = c_2.
    \end{equation*}
    If we take $c_1 = c_2 = 0$, then the problem has infinitely many 
    solutions of the form $u(t) = B \sin(t)$. 
    At the same time, if $c_1 \neq 0$ and $c_2 \neq -c_1$ then 
    the BVP has no solution. To see this observe that the 
    coefficients $A$, $B$ are given by solving the linear 
    system 
    \begin{equation*}
        \begin{bmatrix}
            1 & 0 \\
            -1 & 0
        \end{bmatrix}
        \begin{bmatrix}
            A \\
            B
        \end{bmatrix}
        =
        \begin{bmatrix}
            c_1 \\
            c_2
        \end{bmatrix}
    \end{equation*}

The above example shows the intuitive difference between IVPs 
and BVPs: We think of IVPs as dynamic problems where the solution is 
obtained by starting with initial data and marching forward in time. 
In contrast, BVPs are more like static problems where the solution must satisfy conditions at multiple points, and there is no inherent direction of "marching" in time.

\begin{definition}[IVP]
    \label{def:ivp}
    The general form of an IVP is a system of ODEs of the form 
    \begin{equation*}
        \frac{d\ub}{dt} = \fb(t, \ub), 
        \quad \ub(t_0) = \ub_0,
    \end{equation*}
    where $\ub: [t_0, t_1] \to \R^d$ is the unknown vector-valued 
    solution, $\fb: [t_0, t_1] \times \R^d \to \R^d$ is the right hand side, 
    and $\ub_0 \in \R^d$ is the given initial condition.
\end{definition}
All of our examples from \Cref{ch:Lec1-intro} fall into the IVP category. 
Note that this ODE is only first order, but this model covers 
higher order ODEs via the usual trick of introducing 
auxiliary variables to reduce a high order ODE to a 
first order system. In other words, given 
\begin{equation*}
    \frac{d^n u}{dt^n} = 
    f\left(t, u, \frac{du}{dt}, \ldots, \frac{d^{n-1} u}{dt^{n-1}}\right), 
\end{equation*}
define the auxiliary variables
\begin{equation*}
    u_1 = u, \quad u_2 = \frac{du}{dt}, \quad \ldots, \quad u_n = \frac{d^{n-1} u}{dt^{n-1}}.
\end{equation*}
Then we can rewrite the original $n$-th order ODE as a first order system:
\begin{equation*}
    \frac{d}{dt} \begin{bmatrix} u_1 \\ u_2 \\ \vdots \\ u_n \end{bmatrix}
    = \begin{bmatrix} u_2 \\ u_3 \\ \vdots \\ f(t, u_1, u_2, \ldots, u_n) \end{bmatrix}.
\end{equation*}.

\begin{definition}[BVP]
\label{def:bvp}
    The general form of a BVP is a system of ODEs of the form 
    \begin{equation*}
        \frac{d\ub}{dt} = \fb(t, \ub), 
        \quad \gb(\ub(t_0), \ub(t_1)) = \mathbf{0},
    \end{equation*}
    where $\ub: [t_0, t_1] \to \R^d$ is the unknown vector-valued 
    solution, $\fb: [t_0, t_1] \times \R^d \to \R^d$ is the right hand side, 
    and $\gb: \R^d \times \R^d \to \R^d$ specifies the boundary conditions at $t_0$ and $t_1$.
\end{definition}

\begin{example}\label{ex:poisson-bvp}
    Perhaps the most well-studied example of a BVP is the Poisson's 
    equation 
    \begin{equation*}
        - \frac{d^2 u}{dx^2} = f(x), \quad x \in (0,1), 
        \quad u(0) = u_0, \; u(1) = u_1.
    \end{equation*}
    This is a second order ODE that describes the potential field $u$ 
    caused by a electric charge distribution $f(x)$ with 
    with Dirichlet boundary conditions specified at $x=0$ and $x=1$.
    It also describes many other physical phenomena such as 
    heat conduction in one-dimensional rods and the deflection of elastic 
    beams under load.
\end{example}


% \bibliographystyle{plainnat}
% \bibliography{references}